% C-26 -- conclusions.tex closing-paragraphs merge (reviewer_simulation_trim_audit.md, candidate B)
% Audit copy only -- live file is sections/conclusions.tex.

\revblue{The findings provide a highly efficient baseline demonstrating that lightweight, bio-inspired neural networks ---arbitrating between locomotion modes and reactive behaviors driven by visual and inertial cues--- can manage multimodal gait transitions on embedded hardware, allowing a real robotic platform to navigate, stabilize, and make decisions robustly in environments with obstacles and adverse terrains. The ability to switch locomotion modes based on inertial stability metrics expands the robot's operational range across flat, rough, and inclined surfaces, while the integration of camera, LiDAR, and IMU into low-latency decision loops enabled dynamic goal prioritization, collision avoidance, and timely, collision-free goal approach, validated in both simulation and physical experiments. When scaled onto robust physical platforms, this low-cost computational paradigm illustrates how decentralized, bio-inspired architectures could optimize energy efficiency and behavioral reliability, serving as a stepping stone toward bridging open-loop exploratory prototypes and the rigorous stabilization required by unconstrained, real-world deployment in power-plant monitoring or automated industrial operations. Although the physical prototype displays hardware-level boundaries ---including low-torque actuation and open-loop locomotion profiles--- these results substantiate that complex, multi-stimulus mobility decisions can be executed on such resource-constrained platforms, with the future-work priorities detailed below.}

% Removed (redundant second "in conclusion" paragraph, previously immediately before \section*{Acknowledgements}):
% \revblue{In conclusion, this paper presented a biologically inspired neural controller modeled after
% basal ganglia mechanisms for adaptive gait and mode arbitration on a hybrid wheel-legged quadruped.
% Although the physical prototype displays hardware-level boundaries ---including low-torque actuation
% and open-loop locomotion profiles--- the framework demonstrates that complex, multi-stimulus mobility
% decisions can be executed without heavy computational overhead or expensive training cycles, with the
% future-work priorities detailed above.}
